\documentclass{article}
\usepackage[utf8]{inputenc}
\usepackage{amsmath, amssymb}
\usepackage{geometry}
\usepackage{hyperref}

\geometry{a4paper, margin=1in}

\title{Interactive Solid Mechanics Laboratory \\ Project Report}
\author{Software Engineering Team}
\date{\today}

\begin{document}

\maketitle

\section{Project Overview}
The \textbf{Interactive Solid Mechanics Laboratory} is a desktop software application developed using Python, PyQt6, and PyOpenGL. Its purpose is to provide undergraduate students with a visual, interactive learning environment to understand complex tensor calculus, structural mechanics, and failure theories.

\subsection{Implemented Modules (MVP)}
In the current Minimum Viable Product (MVP), the following modules have been implemented:
\begin{itemize}
    \item \textbf{Mathematical Preliminaries \& Tensors}: Calculation of principal stresses via eigenvalue routines using NumPy.
    \item \textbf{Traction \& Stress State}: A 3D interactive PyOpenGL application to render a stress element.
    \item \textbf{Principal Stresses \& Mohr's Circle}: A 2D Matplotlib canvas dynamically displaying the tri-axial Mohr's Circle state for any generic 3D tensor input.
\end{itemize}

\section{Governing Physics Equations}
This section outlines the fundamental solid mechanics formulas synthesized from the course literature that power the application's calculation engine.

\subsection{Coordinate Transformations}
Vectors and stress tensors rotate according to the orthogonal transformation matrix $\mathbf{Q}$:
\begin{align}
    \mathbf{v}' &= \mathbf{Q}\mathbf{v} \\
    \mathbf{T}' &= \mathbf{Q}\mathbf{T}\mathbf{Q}^T
\end{align}

\subsection{Traction Vector \& Stress Tensor}
The traction vector $\mathbf{t}^{(n)}$ acting on a plane with normal $\mathbf{n}$ is obtained via Cauchy's relationship:
\begin{equation}
    \mathbf{t}^{(n)} = \mathbf{\sigma}^T \mathbf{n}
\end{equation}

\subsection{Stress Equilibrium}
The differential equations of equilibrium, incorporating body forces $\mathbf{b}$:
\begin{equation}
    \nabla \cdot \mathbf{\sigma} + \mathbf{b} = \mathbf{0}
\end{equation}

\subsection{Principal Stresses}
Principal stresses ($\lambda_1, \lambda_2, \lambda_3$) are the roots of the characteristic equation, requiring the determinant of the shifted stress tensor to vanish:
\begin{equation}
    \det(\mathbf{\sigma} - \lambda\mathbf{I}) = 0
\end{equation}

\subsection{Strain Tensor \& Hooke's Law}
The infinitesimal strain tensor relates to the displacement field $\mathbf{u}$ by:
\begin{equation}
    \mathbf{\epsilon} = \frac{1}{2}(\nabla\mathbf{u} + (\nabla\mathbf{u})^T)
\end{equation}
For isotropic, linear elastic materials, generalized Hooke's law links stress and strain:
\begin{equation}
    \mathbf{\epsilon} = \frac{1}{E}[(1+\nu)\mathbf{\sigma} - \nu \text{Tr}(\mathbf{\sigma})\mathbf{I}]
\end{equation}

\subsection{Theory of Beams (Deflection)}
For standard slender beams, the Euler-Bernoulli Theory (EBT) is used:
\begin{equation}
    EI \frac{d^2y}{dX^2} = M(X)
\end{equation}
For deeper beams where shear deformation is significant, the Timoshenko Beam Theory (TBT) is integrated:
\begin{align}
    EI \frac{d\theta}{dX} &= M(X) \\
    \frac{dy}{dX} - \theta &= \frac{V}{k_s G A}
\end{align}

\subsection{Theories of Failure}
To prevent material yielding, stress states must be constrained within failure envelopes.
\\
\textbf{Maximum Shear Stress (Tresca):}
\begin{equation}
    \max \left| \frac{\sigma_i - \sigma_j}{2} \right| < \frac{\sigma_y}{2}
\end{equation}
\textbf{Distortional Energy (von Mises):}
\begin{equation}
    U_d = \frac{1+\nu}{3E}(\sigma_1^2 + \sigma_2^2 + \sigma_3^2 - \sigma_1\sigma_2 - \sigma_2\sigma_3 - \sigma_3\sigma_1) < \frac{1+\nu}{3E}\sigma_y^2
\end{equation}

\end{document}
